% ═══════════════════════════════════════════════════════════════
% MUSTERLÖSUNG: Waagerechter Wurf (Physik Klasse 10)
% ═══════════════════════════════════════════════════════════════

\documentclass[10pt, a4paper]{article}

% ═══════════════════════════════════════════════════════════════
% PAKETE
% ═══════════════════════════════════════════════════════════════

\usepackage[utf8]{inputenc}
\usepackage[T1]{fontenc}
\usepackage[ngerman]{babel}
\usepackage{helvet}
\renewcommand{\familydefault}{\sfdefault}

\usepackage[margin=1.5cm, top=1.8cm, bottom=1.5cm]{geometry}
\usepackage{enumitem}
\usepackage{tabularx}
\usepackage{booktabs}
\usepackage{array}
\usepackage{multicol}
\usepackage{amsmath}
\usepackage{amssymb}

\usepackage{xcolor}
\usepackage{tcolorbox}
\usepackage{fancyhdr}
\usepackage{lastpage}
\usepackage{tikz}

% Kompaktere Listen
\setlist{nosep, leftmargin=1.5em}

% ═══════════════════════════════════════════════════════════════
% FARBEN
% ═══════════════════════════════════════════════════════════════

\definecolor{physik}{HTML}{2c5aa0}
\definecolor{grau}{HTML}{888888}
\definecolor{hellgrau}{HTML}{f0f0f0}
\definecolor{loesung}{HTML}{c62828}
\definecolor{success}{HTML}{2a7a4b}

\newcommand{\fachfarbe}{physik}

% ═══════════════════════════════════════════════════════════════
% VARIABLEN
% ═══════════════════════════════════════════════════════════════

\newcommand{\thema}{Waagerechter Wurf}
\newcommand{\fach}{Physik}
\newcommand{\klasse}{Klasse 10}
\newcommand{\stunde}{Bewegungslehre}

% ═══════════════════════════════════════════════════════════════
% KOPF- UND FUSSZEILE
% ═══════════════════════════════════════════════════════════════

\pagestyle{fancy}
\fancyhf{}
\renewcommand{\headrulewidth}{0.4pt}
\renewcommand{\footrulewidth}{0pt}

\lhead{\small\fach{} \klasse{} -- \thema{} \textcolor{loesung}{\textbf{(MUSTERLÖSUNG)}}}
\rhead{\small}

\cfoot{\small\thepage/\pageref{LastPage}}

% ═══════════════════════════════════════════════════════════════
% BEFEHLE
% ═══════════════════════════════════════════════════════════════

% Lösung (rot)
\newcommand{\lsg}[1]{\textcolor{loesung}{\textbf{#1}}}

% Physik-Highlight
\newcommand{\ph}[1]{\textcolor{\fachfarbe}{\textbf{#1}}}

% Merksatz-Box
\newtcolorbox{merksatz}[1][]{
    colback=hellgrau,
    colframe=\fachfarbe,
    boxrule=0.8pt,
    arc=0pt,
    left=6pt, right=6pt, top=6pt, bottom=6pt,
    fonttitle=\bfseries\small,
    title=#1
}

% Formel-Box
\newtcolorbox{formelbox}{
    colback=white,
    colframe=\fachfarbe,
    boxrule=1.5pt,
    arc=0pt,
    left=8pt, right=8pt, top=8pt, bottom=8pt
}

% Abschnittstitel
\newcommand{\abschnitt}[1]{%
    \vspace{8pt}
    \noindent\textbf{\textcolor{\fachfarbe}{#1}}
    \vspace{4pt}
}

% Punkte am Rand
\newcommand{\punkte}[1]{\hfill\small\textcolor{grau}{/#1 P.}}

% Sternchen für Schwierigkeit
\newcommand{\stmark}{$\star$}

% ═══════════════════════════════════════════════════════════════
% DOKUMENT
% ═══════════════════════════════════════════════════════════════

\begin{document}

% ═══════════════════════════════════════════════════════════════
% TITEL
% ═══════════════════════════════════════════════════════════════

\begin{center}
{\large\bfseries Arbeitsblatt: \thema{} -- \textcolor{loesung}{MUSTERLÖSUNG}}\\[2pt]
{\small\textcolor{grau}{Lösungen in Rot}}
\end{center}

\vspace{-2pt}
\hrule
\vspace{8pt}

% ═══════════════════════════════════════════════════════════════
% ZWEISPALTIG
% ═══════════════════════════════════════════════════════════════

\begin{multicols}{2}
\small

% ─────────────────────────────────────────────────────────────
% KASTEN 1: Die Kernerkenntnis
% ─────────────────────────────────────────────────────────────

\abschnitt{Kasten 1: Das Überlagerungsprinzip}

\begin{merksatz}[Merksatz]
Beim waagerechten Wurf überlagern sich zwei \ph{unabhängige} Bewegungen:

\vspace{6pt}
\begin{tabular}{ll}
\textbf{Horizontal:} & \lsg{gleichförmige} Bewegung \\[4pt]
\textbf{Vertikal:} & \lsg{freier Fall} \\
\end{tabular}

\vspace{8pt}
Die horizontale Geschwindigkeit beeinflusst die

\lsg{Fallzeit} \textbf{nicht}.

\vspace{6pt}
Deshalb landen ein fallender und ein geworfener Ball,

die gleichzeitig starten, \lsg{gleichzeitig}.
\end{merksatz}

\vspace{4pt}

% ─────────────────────────────────────────────────────────────
% SKIZZE
% ─────────────────────────────────────────────────────────────

\abschnitt{Skizze: Flugbahn}

\begin{center}
\begin{tikzpicture}[scale=0.65]
    % Koordinatensystem
    \draw[->] (0,0) -- (6.5,0) node[right] {$x$};
    \draw[->] (0,0) -- (0,4.5) node[above] {$y$};

    % Boden
    \draw[thick] (-0.5,0) -- (6.5,0);

    % Plattform
    \fill[gray!50] (-0.5,4) rectangle (0.5,4.2);

    % Höhe h
    \draw[<->] (-0.3,0) -- (-0.3,4) node[midway, left] {$h$};

    % Parabel
    \draw[thick, loesung] (0.5,4) parabola (5.5,0);

    % Ball Start
    \fill[physik] (0.5,4) circle (0.15);

    % Ball Ende
    \fill[physik] (5.5,0) circle (0.15);

    % Wurfweite
    \draw[<->] (0.5,-0.5) -- (5.5,-0.5) node[midway, below] {$s_x$};

    % v0 Pfeil (horizontal)
    \draw[->, thick, loesung] (0.5,4) -- (2,4) node[above] {$\vec{v}_x$};

    % vy Pfeil (vertikal, an Zwischenposition)
    \draw[->, thick, blue] (3,2.2) -- (3,0.7) node[right] {$\vec{v}_y$};

    % vx Pfeil an Zwischenposition
    \draw[->, thick, loesung] (3,2.2) -- (4.5,2.2) node[above] {$\vec{v}_x$};
\end{tikzpicture}
\end{center}

% ─────────────────────────────────────────────────────────────
% KASTEN 2: Formeln
% ─────────────────────────────────────────────────────────────

\abschnitt{Kasten 2: Die Formeln}

\begin{formelbox}
\textbf{Schritt 1: Fallzeit berechnen}
\vspace{4pt}

\begin{equation*}
t = \sqrt{\frac{2h}{g}}
\end{equation*}

\vspace{2pt}
\small
$t$ = Fallzeit in s \quad $h$ = Höhe in m \quad $g = 10\,\text{m/s}^2$
\end{formelbox}

\vspace{6pt}

\begin{formelbox}
\textbf{Schritt 2: Wurfweite berechnen}
\vspace{4pt}

\begin{equation*}
s_x = v_0 \cdot t
\end{equation*}

\vspace{2pt}
\small
$s_x$ = Wurfweite in m \quad $v_0$ = Anfangsgeschw. in m/s
\end{formelbox}

\vspace{6pt}

% ─────────────────────────────────────────────────────────────
% WICHTIGE WERTE
% ─────────────────────────────────────────────────────────────

\abschnitt{Tabelle: Typische Fallzeiten (mit $g = 10\,\text{m/s}^2$)}

\begin{center}
\begin{tabular}{|c|c|}
\hline
\textbf{Höhe $h$} & \textbf{Fallzeit $t$} \\
\hline
5 m & 1 s \\
\hline
20 m & 2 s \\
\hline
45 m & 3 s \\
\hline
80 m & 4 s \\
\hline
\end{tabular}
\end{center}

\vspace{4pt}

% ─────────────────────────────────────────────────────────────
% AUFGABEN MIT LÖSUNGEN
% ─────────────────────────────────────────────────────────────

\columnbreak

\abschnitt{Aufgabe 1: Grundaufgabe (\stmark\stmark)}\punkte{4}

Ein Ball wird aus $h = 20\,\text{m}$ Höhe waagerecht mit $v_0 = 15\,\text{m/s}$ geworfen.

\textbf{a)} Berechne die Fallzeit.

\vspace{4pt}
\textcolor{loesung}{
$t = \sqrt{\dfrac{2h}{g}} = \sqrt{\dfrac{2 \cdot 20\,\text{m}}{10\,\text{m/s}^2}}$

$t = \sqrt{\dfrac{40}{10}\,\text{s}^2} = \sqrt{4\,\text{s}^2}$

$\boxed{t = 2\,\text{s}}$
}

\vspace{6pt}

\textbf{b)} Berechne die Wurfweite.

\vspace{4pt}
\textcolor{loesung}{
$s_x = v_0 \cdot t = 15\,\text{m/s} \cdot 2\,\text{s}$

$\boxed{s_x = 30\,\text{m}}$
}

\vspace{8pt}

% ─────────────────────────────────────────────────────────────

\abschnitt{Aufgabe 2: Anwendung (\stmark\stmark)}\punkte{5}

Ein Flugzeug fliegt in $h = 45\,\text{m}$ Höhe mit einer Geschwindigkeit von $v_0 = 20\,\text{m/s}$. Es lässt ein Hilfspaket fallen.

\textbf{a)} Wie lange ist das Paket in der Luft?

\vspace{4pt}
\textcolor{loesung}{
$t = \sqrt{\dfrac{2 \cdot 45\,\text{m}}{10\,\text{m/s}^2}} = \sqrt{9\,\text{s}^2}$

$\boxed{t = 3\,\text{s}}$
}

\vspace{6pt}

\textbf{b)} Wie weit vor dem Ziel muss das Flugzeug das Paket abwerfen?

\vspace{4pt}
\textcolor{loesung}{
$s_x = v_0 \cdot t = 20\,\text{m/s} \cdot 3\,\text{s}$

$\boxed{s_x = 60\,\text{m}}$
}

\vspace{8pt}

% ─────────────────────────────────────────────────────────────

\abschnitt{Aufgabe 3: Transferaufgabe (\stmark\stmark\stmark)}\punkte{6}

Ein Stuntman springt von einer Klippe ($h = 80\,\text{m}$) waagerecht ab. Er landet im Meer, 120 m von der Klippenkante entfernt.

\textbf{a)} Berechne die Fallzeit.

\vspace{4pt}
\textcolor{loesung}{
$t = \sqrt{\dfrac{2 \cdot 80\,\text{m}}{10\,\text{m/s}^2}} = \sqrt{16\,\text{s}^2}$

$\boxed{t = 4\,\text{s}}$
}

\vspace{6pt}

\textbf{b)} Berechne die Absprunggeschwindigkeit $v_0$.

\vspace{4pt}
\textcolor{loesung}{
Umstellen: $s_x = v_0 \cdot t \Rightarrow v_0 = \dfrac{s_x}{t}$

$v_0 = \dfrac{120\,\text{m}}{4\,\text{s}}$

$\boxed{v_0 = 30\,\text{m/s}}$
}

\vspace{8pt}

% ─────────────────────────────────────────────────────────────

\abschnitt{Selbstcheck}

\begin{tabular}{|p{5.5cm}|c|c|}
\hline
\textbf{Ich kann...} & \textbf{Ja} & \textbf{Nein} \\
\hline
...erklären, warum beide Bälle gleichzeitig landen. & \lsg{$\boxtimes$} & $\square$ \\
\hline
...die Fallzeit aus der Höhe berechnen. & \lsg{$\boxtimes$} & $\square$ \\
\hline
...die Wurfweite berechnen. & \lsg{$\boxtimes$} & $\square$ \\
\hline
\end{tabular}

\end{multicols}

% ═══════════════════════════════════════════════════════════════
% PUNKTE-ÜBERSICHT (Fußbereich)
% ═══════════════════════════════════════════════════════════════

\vfill
\hrule
\vspace{4pt}
\small
\begin{tabular}{llll}
\textbf{Gesamt:} & \lsg{15} / 15 P. &
\textbf{Note/NP:} & \lsg{15 NP / Note 1}
\end{tabular}

\end{document}
